\section*{Chapter \ref{ch:b3:radicals-carbenes} --- Radicals, Carbenes and Rearrangements}

\begin{solution}{exo:b3:radicals-carbenes:1}
Propane has six primary and two secondary hydrogens. Chlorination: $6 \times 1 : 2 \times 3.9 = 6 : 7.8$,
that is 43\,\% 1-chloropropane and 57\,\% 2-chloropropane. Bromination: $6 : 164$, 4\,\% and
96\,\%.
\end{solution}

\begin{solution}{exo:b3:radicals-carbenes:2}
$\ce{CH2}$: triplet; $\ce{CCl2}$: singlet (the chlorine lone pairs raise the empty p orbital).
Dichlorocarbene adds stereospecifically: \textit{trans}-1,1-dichloro-2,3-dimethylcyclopropane.
\end{solution}

\begin{solution}{exo:b3:radicals-carbenes:3}
\textit{cis}-1,2-Diethylcyclopropane: the \omterm{def:b3:radicals-carbenes:carbenoid}{carbenoid} adds in one step to one face.
\end{solution}

\begin{solution}{exo:b3:radicals-carbenes:4}
Cyclohexanone: $\varepsilon$-caprolactone (oxepan-2-one). Acetophenone: phenyl acetate,
$\ce{CH3CO2C6H5}$; phenyl migrates in preference to methyl.
\end{solution}

\begin{solution}{exo:b3:radicals-carbenes:5}
Nine primary and one tertiary hydrogen. Chlorination: $9 : 5.1$, 64\,\% 1-chloro-2-methylpropane
and 36\,\% 2-chloro-2-methylpropane. Bromination: $9 : 1600$, 0.6\,\% and 99.4\,\%.
\end{solution}

\begin{solution}{exo:b3:radicals-carbenes:6}
$k_c = 3.0 \times k_H[\mathrm{Bu_3SnH}] = 3.0 \times \num{2.4e6} \times 0.020 = \qty{1.4e5}{s^{-1}}$.
\end{solution}

\begin{solution}{exo:b3:radicals-carbenes:7}
Attack on C5: 5-exo-trig, favoured; on C6: 6-endo-trig, favoured by the rules but
slower here. 4-exo-tet: favoured. 5-endo-trig: disfavoured. 5-exo-dig: favoured.
3-exo-dig: disfavoured. 5-endo-dig: favoured.
\end{solution}

\begin{solution}{exo:b3:radicals-carbenes:8}
Loss of water gives a tertiary, benzylic cation; on the neighbouring carbon a phenyl
and a methyl group could move, and phenyl, of higher aptitude, migrates. The
oxygen-stabilised cation loses a proton: 3,3-diphenylbutan-2-one.
\end{solution}

\begin{solution}{exo:b3:radicals-carbenes:9}
Curtius: $\ce{C6H5CON3 -> C6H5NCO + N2}$. With water, the carbamic acid loses $\ce{CO2}$: aniline
(some diphenylurea forms as aniline meets more \omterm{def:b3:radicals-carbenes:curtius}{isocyanate}). With ethanol: ethyl
\textit{N}-phenylcarbamate, $\ce{C6H5NHCO2C2H5}$.
\end{solution}

\begin{solution}{exo:b3:radicals-carbenes:10}
Chlorine: primary $421.7 - 432 = \qty{-10}{kJ/mol}$, tertiary $405.4 - 432 = \qty{-27}{kJ/mol}$; bromine: $+56$ and
$+39$. Both chlorine abstractions are exothermic, with early transition states that
feel little of the \qty{16}{kJ/mol} difference; both bromine abstractions are endothermic,
with late transition states whose activation energies differ by most of it:
$\eu^{16000/(8.314 \times 298)} \approx 600$, the order of the observed ratio.
\end{solution}

\begin{solution}{exo:b3:radicals-carbenes:11}
\textit{E} oxime (OH anti to C2, which bears the methyl group): C2 migrates, nitrogen enters
next to it: 7-methylazepan-2-one. \textit{Z} oxime: C6 migrates: 3-methylazepan-2-one. In the
\omterm{def:b3:radicals-carbenes:beckmann}{Beckmann rearrangement} the oxime geometry decides; in the \omterm{def:b3:radicals-carbenes:baeyer}{Baeyer--Villiger
oxidation} the \omterm{def:b3:radicals-carbenes:migratory}{migratory aptitude} does, and the secondary carbon moves:
7-methyloxepan-2-one.
\end{solution}

\begin{solution}{exo:b3:radicals-carbenes:12}
Fraction cyclised $k_c/(k_c + k_H c)$: 0.91, 0.49 and 0.09. For 95\,\%: $k_Hc = k_c(1/0.95 - 1)$, so
$c = 0.0526 \times \num{2.3e5}/\num{2.4e6} = \qty{5.0e-3}{mol/L}$. It is maintained by adding the hydride slowly with a
syringe pump, or by using a catalytic amount of tin halide with a reducing agent
that regenerates the hydride.
\end{solution}

\begin{solution}{pb:b3:radicals-carbenes:1}
\textbf{1.} $\ce{C6H10O + NH2OH -> C6H11NO + H2O}$.

\textbf{2.} Acid is needed to remove the hydroxy group of the addition intermediate as
water, but too much acid protonates hydroxylamine, which then no longer adds.

\textbf{3.} No: the two carbons attached to the C=N carbon are equivalent by the
symmetry of the ring.

\textbf{4.} In the \textit{E} oxime the OH points away from the methyl-bearing carbon C2 (anti to
it), in the \textit{Z} oxime towards it.

\textbf{5.} Inversion at an oxime nitrogen bearing an electronegative oxygen has a high
barrier.

\textbf{6.} It protonates (or sulfonates) the hydroxy group, making it a good leaving
group, water.

\textbf{7.} The ring C--C bond anti to the leaving group moves from carbon to nitrogen,
which is thereby inserted into the ring: six atoms become seven.

\textbf{8.} A cyclic nitrilium ion.

\textbf{9.} $\ce{C6H11NO -> C6H11NO}$: an isomerisation, oxime to lactam.

\textbf{10.} Caprolactam, azepan-2-one, the seven-membered cyclic amide of
6-aminohexanoic acid; its ring-opening polymerisation gives nylon-6.

\textbf{11.} \qty{98.0}{g/mol} and \qty{113.0}{g/mol}.

\textbf{12.} $\qty{1.0e6}{g}/\qty{98.0}{g/mol} = \qty{10204}{mol}$; $\times 0.98 \times 0.98 = \qty{9800}{mol}$; $\times \qty{113.0}{g/mol} = \qty{1.11e6}{g}$, \qty{1.11}{t}.

\textbf{13.} Oxime: $\num{10204} \times 0.98 = \qty{1.00e4}{mol}$; $\ce{H2SO4}$: $\qty{1.30e4}{mol}$, giving as many moles of
$\ce{(NH4)2SO4}$: $\qty{1.30e4}{mol} \times \qty{132.1}{g/mol} = \qty{1.7}{t}$, more than the caprolactam.

\textbf{14.} $\ce{C6H10O + RCO3H -> C6H10O2 + RCOOH}$.

\textbf{15.} The tetrahedral adduct of the peroxy acid on the carbonyl carbon (the
Criegee intermediate), a peroxy hemiketal.

\textbf{16.} A ring $\ce{CH2}$ group (the two are equivalent) moves to the nearer peroxide
oxygen; the carboxylate $\ce{RCOO-}$ leaves.

\textbf{17.} $\varepsilon$-Caprolactone, oxepan-2-one; its ring-opening polymerisation gives
poly($\varepsilon$-caprolactone), a biodegradable polyester.

\textbf{18.} Its only by-product is water, where a peroxy acid leaves a mole of
carboxylic acid per mole of lactone and is itself hazardous to store.

\textbf{19.} The secondary carbon C2 migrates: 7-methyloxepan-2-one.

\textbf{20.} Yes: the migrating carbon keeps its configuration.

\textbf{21.} 7-Methylazepan-2-one (C2 anti, migrates).

\textbf{22.} 3-Methylazepan-2-one (C6 migrates).

\textbf{23.} Beckmann: the geometry of the oxime (the anti group moves).
Baeyer--Villiger: the \omterm{def:b3:radicals-carbenes:migratory}{migratory aptitude} (the more substituted carbon moves).

\textbf{24.} Its two $\alpha$ carbons are equivalent: one oxime, one lactam, one lactone.

\textbf{25.} One tonne of cyclohexanone gives about \qty{1.11}{t} of caprolactam.
\end{solution}
